\chapter{The original synthetic teaching gallery}\label{ch:data}

\begin{notebox}
This chapter describes the retained 57-run gallery and its 325-row invented history, separate from the new .ixo-calibrated demo in Chapter~\ref{ch:roche}. Everything in this chapter is \textbf{made-up data for demonstrative purposes}.
\end{notebox}

\section{Why synthetic data, and how it resembles the real files}
Real experiment files carry sample names, laboratory identities, instrument serial numbers and results that must not be
published. The documentation uses synthetic raw runs calibrated to the measured scale, moving-range noise and resolution
of two supplied .ixo instrument summaries. No real values or dates are copied. No .eds instrument history was supplied, so
.eds instrument keeps the former demonstration scale. All values, dates, identifiers, curves and events shown here are invented.
The data are built in three steps:

\begin{enumerate}
\item \textbf{Structure templates.} Five experiment files (two .ixo instrument, three .eds/.edt instrument) were decoded by the page
      once to obtain the \emph{shape} of a run: plate format (8$\times$12), thermal program, detection channels
      (e.g.\ 465-510 / 533-580 on the .ixo instrument, x1-m1 \dots\ x5-m5 on the .eds instrument), the analyses and the fields each file type
      carries. Seeded rules, the .ixo instrument calibration and logistic curves generate the displayed measurements.
      .eds instrument Cq centres are independently defined by synthetic rules. Missing template programs receive an explicitly synthetic 40-cycle program. Thus these are synthetic demonstrations, not
      independent experimental evidence or a claim of statistical anonymisation.
\item \textbf{Regeneration.} Each synthetic run is a copy of one template in which \emph{every} value that appears in this document
      is replaced (Listings below):
      \begin{itemize}
      \item names: samples become \texttt{DEMO-L1-001} \dots, controls become \texttt{NTC}, \texttt{Blank extraction},
            \texttt{DEMO LOD 01.03.2026} or \texttt{Positive control}; operators become \emph{Analyst A/B/C};
            instruments become \texttt{LC-DEMO-01}, \texttt{LC-DEMO-02}, \texttt{QS-DEMO-A}, \texttt{QS-DEMO-B}, \texttt{QS-DEMO-C};
            all UIDs, plate IDs, hashes, paths and host names are random;
      \item dates: all timestamps are moved to the synthetic run date, keeping their order and spacing;
      \item amplification curves: logistic curves with the chosen Cq at the second-derivative maximum, baseline, amplitude and noise
            drawn at random; the stored Cq is set to the generated Cq, and for .eds instrument the stored Ct is the threshold crossing
            of the generated $\Delta$Rn;
      \item .ixo instrument background, plateau and exposure use calibrated channel scales and a shared synthetic plate-load
            factor. The thermal model uses calibrated heat and cooling rates, overshoot, settling, hold offset and transitions;
            it produces new temperature traces and optics measurements. .eds instrument retains its earlier scale;
      \item instrument logs, lamp reference, exposure, LED, filter wheel, cover, timing: regenerated around typical values;
      \item stored relative quantification: recomputed from the synthetic Cts ($2^{-\Delta\Delta C_q}$).
      \end{itemize}
      A leak check then searches every synthetic run for the identifiers of the templates; the data set is used only when
      nothing is found.
\item \textbf{History.} A multi-year instrument history (325 rows) is generated in Python with the same .ixo instrument
      calibration and the former .eds instrument model, then given deliberate trends, a service step, seasonality and operator differences
      (Listing \texttt{lc1()} below). It is imported into the page exactly as a laboratory would import its own history file.
\end{enumerate}

{\footnotesize
\begin{longtable}{@{}llllll@{}}
\toprule \textbf{Column 1} & \textbf{Column 2} & \textbf{Column 3} & \textbf{Column 4} & \textbf{Column 5} & \textbf{Column 6} \\\midrule\endfirsthead
\toprule \textbf{Column 1} & \textbf{Column 2} & \textbf{Column 3} & \textbf{Column 4} & \textbf{Column 5} & \textbf{Column 6} \\\midrule\endhead
LC · LC-DEMO-01 & 30 & 150 & 2024-01-08 & 2026-08-29 & Analyst A, Analyst B, Analyst C \\\hline
LC · LC-DEMO-02 & 8 & 55 & 2025-01-14 & 2026-08-10 & Analyst A, Analyst B, Analyst C \\\hline
QS · QS-DEMO-A & 10 & 90 & 2024-03-08 & 2026-08-28 & Analyst A, Analyst B, Analyst C \\\hline
QS · QS-DEMO-B & 6 & 30 & 2025-06-06 & 2026-07-31 & Analyst A, Analyst B, Analyst C \\\hline
QS · QS-DEMO-C & 3 & 0 & 2026-07-06 & 2026-08-17 & Analyst A, Analyst B, Analyst C \\\hline
\bottomrule
\end{longtable}}



\section{Planted events}
So that the charts and the forensic checks have something to find, the following events were written into the synthetic data:

\begin{center}\footnotesize
\begin{tabular}{@{}p{3.3cm}p{5.6cm}p{6.6cm}@{}}\toprule
Where & Event & Where it shows\\\midrule
LC-DEMO-01, 2025-05-12 & synthetic service step: settling $4.95\to4.54$\,s, heating rate $4.89\to4.77$\,°C/s, hold offset $-0.50\to-0.40$\,°C & B-H1, B-H2, B-H3b, B-H4 (Figs.~\ref{fig:cool_before}, \ref{fig:cool_rebase}, \ref{fig:hold})\\
LC-DEMO-01, whole period & calibrated cooling rate declines from about 2.47 to 2.22\,°C/s; 95-to-60\,°C transition is about 15.4\,s & B-H2 trend and forecast, B-H3c transition times\\
LC-DEMO-01, loaded runs & two extraction batches with simulated age-related Cq drift & Control batch and extract-age graphs\\
Analyst C & NTC positive rate 7\,\% (others 0.8\,\%) & B-O5 funnel, B-M5 p-chart\\
Analyst B & replicate SD varied around the calibrated 0.18--0.20 scale & B-O1\\
\texttt{DEMO\_LC1\_SCREEN\_013} & well A03 carries a copy of the curve of A01 & forensic log: \emph{Copied data} (error), B-O11\\
\texttt{DEMO\_LC1\_SCREEN\_018} & last change 410 h after the run & forensic log (timeline), A-S4\\
\texttt{DEMO\_LC1\_SCREEN\_022} & extraction blank amplifies at Cq 35.6 & SOP: review required; forensic finding\\
QS-DEMO-A & LED current rising, cover drive slowing, filter wheel +14\,ms from 2025-09 & B-H10, B-H18, B-H14\\
QS-DEMO-A, run 8 & three command errors in the instrument log & B-S8\\\bottomrule
\end{tabular}
\end{center}

\section{Generator code}
The core functions of the generator follow; the complete files are in Appendix~\ref{app:generator}.

\par\medskip\noindent{\small\textbf{Listing}~\texttt{curve()}\hfill\textcolor{gray}{demo\_synth.js lines 38--41}}\par\nopagebreak
\lstinputlisting[style=vstyle,language=JavaScript,firstnumber=38]{code/demo_curve.js}


The logistic curve $F(c)=B+\dfrac{A}{1+e^{-(c-x_0)/s}}$ has its second-derivative maximum at $c=x_0-1.317\,s$, so
$x_0=C_q+1.317\,s$ places the .ixo instrument ``2nd derivative maximum'' crossing point exactly at the chosen $C_q$.

\par\medskip\noindent{\small\textbf{Listing}~\texttt{thermalTrace()}\hfill\textcolor{gray}{demo\_synth.js lines 43--59}}\par\nopagebreak
\lstinputlisting[style=vstyle,language=JavaScript,firstnumber=43]{code/demo_thermalTrace.js}


\par\medskip\noindent{\small\textbf{Listing}~\texttt{lcRun()}\hfill\textcolor{gray}{demo\_synth.js lines 61--141}}\par\nopagebreak
\lstinputlisting[style=vstyle,language=JavaScript,firstnumber=61]{code/demo_lcRun.js}


\par\medskip\noindent{\small\textbf{Listing}~\texttt{lc1() — model of the synthetic history of LC-DEMO-01}\hfill\textcolor{gray}{make\_history.py lines 36--54}}\par\nopagebreak
\lstinputlisting[style=vstyle,language=python,firstnumber=36]{code/hist_lc1.py}



\fig{g_all_x_duration.png}{Run duration and instrument use of the 57 synthetic runs, coloured by instrument (Graphs tab, \texttt{x\_duration}).}{dataset_duration}


